\loai{Cuộn dây có 1 vòng}
\begin{vidu} %[2P4-2.3B][Loại1] 
Tại tâm của một dòng điện tròn cường độ 5 A người ta đo được cảm ứng từ $B = 31,4.10^{- 6}$ T. Đường kính của dòng điện tròn là
\choice
{10 cm}
{2 cm}
{1 cm}
{\True 20 cm}
\loigiai{Ta có: $B = 2\pi.10^{-7}.\dfrac{I}{R} \Rightarrow R = 2\pi.10^{-7}.\dfrac{I}{B} = 2\pi.10^{-7}.\dfrac{5}{31,4.10^{-6}} = 0,1\ m = 10\ cm.$\\
$\Rightarrow$ đường kính dòng điện $d = 2R = 20$ cm}
\end{vidu}
\loai{Cuộn dây có nhiều vòng}
\begin{vidu} %[2P4-2.3B][Loại2]
Một dòng điện chạy trong một dây tròn 20 vòng đường kính 20 cm với cường độ 10 A thì cảm ứng từ tại tâm các vòng dây là
\choice
{\True $0,4\pi\ mT$}
{$0,02\pi\ mT$}
{$20\pi\ \mu\ T$}
{0,2 mT}
\loigiai{Cảm ứng từ tại tâm các vòng dây là: $B=2\pi.10^{-7}N\dfrac{I}{R}=0,4\pi.10^{-3}\ T$}
\end{vidu}

\begin{vidu} %[2P4-2.3B][Loại2]
Tại tâm của một khung dây tròn gồm 100 vòng (quấn sít nhau) đo được cảm ứng từ có độ lớn $6,28.10^{-5}$ T. Đường kính của khung dây tròn là 10 cm. Lấy $\pi = 3,14$. Cường độ dòng điện chạy qua mỗi vòng là
\choice
{5 A}
{1 A}
{\True 0,05 A}
{0,5 A}
\loigiai{$R = 5\ cm,\ B=2\pi .10^{-7}N\dfrac{I}{R}\Rightarrow I= 0,05$ A}
\end{vidu}
\loai{Mối liên hệ giữa cảm ứng từ với cường độ dòng điện}
\begin{vidu} %[2P4-2.3B][Loại3]
Một dây dẫn tròn mang dòng điện 20 A thì tâm vòng dây có cảm ứng từ $0,4\pi\ \mu T$. Nếu dòng điện qua giảm 5 A so với ban đầu thì cảm ứng từ tại tâm vòng dây là
\choice
{\True $0,3\pi \ \mu T$}
{$0,5\pi \ \mu T$}
{$0,2\pi \ \mu T$}
{$0,6\pi \ \mu T$}
\loigiai{Ta có: $B=2\pi.10^{-7}.\dfrac{I}{R}\Rightarrow \dfrac{B_1}{B_2}=\dfrac{I_1}{I_2}\Leftrightarrow \dfrac{15}{20}=\dfrac{B_2}{0,4\pi}\Rightarrow B_2=0,3\pi \ \mu T$}
\end{vidu}

\loai{Cuộn dây phải tính số vòng}
\begin{vidu} %[2P4-2.3K][Loại4]
Một dây dẫn (tiết diện dây dẫn rất nhỏ) chiều dài 18,84 m được bọc bằng một lớp cách điện mỏng, quấn tròn thành một cuộn dây có bán kính 10 cm. Cho dòng điện có cường độ 0,4 A đi qua dây. Lấy $\pi = 3,14$. Cảm ứng từ tại tâm cuộn dây là
\choice
{\True $7,54.10^{-5}$ T}
{$7,54.10^{-4}$ T}
{$4,57.10^{-5}$ T}
{$4,57.10^{-4}$ T}
\loigiai{\begin{itemize}
\item Chu vi 1 vòng dây là $P = 2\pi R = 0,628\ m$.
\item Số vòng được quấn là $N=\dfrac{L}{P}=30$ vòng.
\item $B=2\pi .10^{-7}N\dfrac{I}{R}=7,{{54.10}^{-5}}$ T.
\end{itemize}}
\end{vidu}


\loai{Cuộn dây có 1 số vòng dây bị cuốn ngược}
\begin{vidu} %[2P4-2.3K][Loại5]
Một khung dây tròn gồm 24 vòng dây có dòng điện chạy qua. Theo tính toán thì cảm ứng từ ở tâm khung bằng $6,3.10^{-5}$ T. Nhưng khi đo thì thấy cảm ứng từ ở tâm bằng $4,2.10^{-5}$ T, kiểm tra lại thấy có một số vòng dây bị quấn nhầm chiều ngược chiều với đa số các vòng trong khung. Số vòng dây bị quấn nhầm là
\choice
{6 vòng}
{2 vòng}
{8 vòng}
{\True 4 vòng}
\loigiai{\begin{itemize}
\item Theo tính toán: $B=2\pi .10^{-7}. 24. \dfrac{I}{R}=6,3.10^{-5}$ T.
\item Quấn ngược n vòng nên: $B=2\pi .10^{-7}\left( 24-2n \right)\dfrac{I}{R}=4,2.10^{-5}\ T. \Rightarrow \dfrac{24}{24-2n}=\dfrac{3}{2}\Rightarrow n=4$.
\end{itemize}}
\end{vidu}
